\section{Discussion}
\label{sec:discussion-jv}
The persistent negative gamma regime in our 2024 sample deserves emphasis. Historical options literature assumed regime-switching dynamics \cite{ni2005stock,garleanu2009demand}, where markets alternated between dealer long gamma (dampening volatility) and dealer short gamma (amplifying volatility). The 0DTE options explosion has fundamentally altered this dynamic \cite{cboe2023}, creating an environment where dealers remain structurally short gamma. Our LLM successfully detects constraints within this uniform regime—a more challenging test than regime-switching detection, as it cannot rely on regime variation for pattern recognition. The combination of high detection rate (69.4\%), significant predictive power (Granger p = 0.042), and 91.2\% materialization accuracy demonstrates a genuine understanding of structural mechanisms rather than statistical pattern-matching.

\subsection{Interpretation of Results}

Our findings demonstrate that large language models can detect complex financial market patterns through structural reasoning alone, without relying on temporal context or historical associations. The 71.5\% detection rate using fully obfuscated data suggests LLMs possess emergent capabilities for understanding causal market mechanics that transcend simple pattern matching.

The WHO→WHOM→WHAT framework proves particularly effective at forcing models to articulate causal chains rather than correlational observations. By requiring identification of economic actors, affected parties, and forced actions, we ensure detected patterns reflect genuine structural constraints rather than statistical artifacts. The consistent 91.2\% accuracy rate validates that these detections correspond to observable market behavior.

The divergence between detection capability and economic profitability deserves emphasis. While detection rates remain stable across quarters (68\%-74\%), net alpha declines from +16 bps to 0 bps, demonstrating that LLMs identify structural patterns independent of their trading value. This separation validates our methodology—we test understanding of market mechanics, not the ability to generate profits.

\subsection{Limitations}
Several limitations warrant consideration:
\textbf{Single Asset Focus:} Testing exclusively on SPY options limits generalizability to other underlyings and asset classes.
\textbf{End-of-Day Granularity:} Daily snapshots may miss intraday dynamics critical to 0DTE hedging patterns. However, our statistical validation confirms EOD data contains sufficient predictive information (AUC = 0.681) for T+1 full-day outcomes, even though overnight gap prediction is weak (r = -0.002). This suggests dealer constraints operate contemporaneously rather than persisting overnight, with the full-day materialization capturing constraint effects that accumulate throughout the trading session.
\textbf{Simplified Gamma Calculations:} Standard Black-Scholes Greeks without smile adjustments may introduce noise in absolute GEX values.
\textbf{LLM Black Box Nature:} The 28.5 percentage point (pp) gap between biased and unbiased prompts suggests significant sensitivity to prompt framing.
\textbf{Temporal Limitation:} Our 2024 sample may not generalize to crisis periods or structural bull/bear markets.
\textbf{LLM vs Statistical Model Trade-off:} While LLMs excel at pattern identification and causal interpretation (91.2\% accuracy), statistical models outperform for probability calibration (AUC 0.681 vs 0.488). This suggests complementary rather than competitive use cases: LLMs for mechanism detection, statistical models for quantitative prediction.

\subsection{Prompting Fidelity in Financial LLM Research}

Our analysis uncovered an important methodological limitation regarding prompt design for financial pattern detection. To test whether the LLM's reasoning qualitatively adapts across different market conditions, we conducted a supplementary experiment requesting detailed 50-100 word causal explanations using prompts that explicitly listed expected qualitative keywords (``amplification'', ``cascading'', ``dampening'').

The result was a null finding: responses for high-alpha quarters (Q1, Sharpe 1.8) and zero-alpha quarters (Q4, Sharpe 0.1) became nearly identical, with both producing template-like responses containing 186-197 amplification keywords and uniform ``strong amplification'' intensity characterizations. Chi-squared tests confirmed no significant differentiation ($p > 0.05$).

This contrasts sharply with the subtle but genuine linguistic adaptation observed in unprompted brief responses (``Forced to buy'' $\rightarrow$ ``Maintain equilibrium''). We interpret this discrepancy as \textit{prompt bias artifact}: explicitly requesting rich qualitative language triggers the LLM's pre-trained knowledge base regarding negative GEX regimes, overriding contextual adaptation signals. The prompt effectively ``teaches to the test'', eliciting learned templates rather than genuine reasoning.

\textbf{Contribution to FinLLM Research.} This finding establishes a new guardrail for financial LLM methodology: \textit{unbiased brief outputs often provide higher signal fidelity than explicitly structured rich outputs}. Researchers designing LLM-based financial analysis systems should prioritize minimal prompting for core reasoning tasks, reserving detailed structured outputs for post-detection explanation or reporting phases. Prompts that enumerate expected keywords or explicitly guide qualitative language risk corrupting the very reasoning signals they aim to elicit.

This methodological contribution extends beyond our specific 0DTE hedging application, offering guidance for any financial LLM research seeking to validate genuine reasoning versus pattern memorization.

\subsection{Implications for Financial Markets}
The ability to detect dealer constraints without historical context suggests LLMs could identify novel patterns in unexplored markets through pure structural reasoning. These constraints emerge from countless independent dealers responding to local gamma exposures rather than centralized coordination—precisely the type of spontaneous market order that resists simple memorization. Risk managers could use LLM-detected patterns as early warning signals for volatility amplification, while the obfuscation testing methodology validates a new approach for distinguishing genuine structural constraints from historical correlations.

As LLMs increasingly influence trading decisions, regulators must understand their pattern detection capabilities. Our framework provides tools for assessing whether AI systems genuinely understand market dynamics or merely memorize historical patterns.

\subsection{Theoretical Contributions}

We introduce \textit{obfuscation testing} as a rigorous methodology for validating structural pattern detection by removing temporal context while preserving mechanical relationships. Our results demonstrate that transformer architectures exhibit emergent understanding of complex financial constraints, detecting dealer hedging patterns from pure gamma exposure values. The WHO→WHOM→WHAT framework proves effective for forcing explicit causal attribution in financial applications. Together, these contributions distinguish genuine causal reasoning from pattern memorization in LLM financial analysis.

Our findings raise interesting questions about the nature of market understanding. While LLMs detect structural patterns algorithmically, actual market participants—dealers and traders—engage in purposeful action under constraints \cite{mises1949,rothbard1962}. The divergence between pattern detection capability and economic profitability (Figure~\ref{fig:quarterly_stability-jv}) suggests LLMs identify what Austrian economists termed ``market regularities''—emergent coordination patterns arising from dispersed knowledge \cite{hayek1945}—without engaging in the entrepreneurial discovery process that generates alpha \cite{mises1949,rothbard1962}. This distinction has implications for deploying AI systems in financial markets: structural pattern recognition, while valuable for risk management and market surveillance, differs fundamentally from the adaptive decision-making that produces trading edge.

\subsection{Future Directions}
Future work should validate detection across multiple asset classes and market regimes to establish generalizability. Intraday analysis with high-frequency data could reveal microstructure patterns invisible at daily granularity, particularly for 0DTE options. Testing ensemble methods across multiple LLMs could distinguish model-specific artifacts from robust pattern detection through consensus mechanisms.
